\Needspace{15\baselineskip}
\section{Power-Shortfall Diagnosis}
\label{sec:disc-plan-gap}

The plan's \SI{15}{\watt} shortfall uses the reported approximately
\SI{85}{\watt} V2 output against the \SI{100}{\watt} target
\cite{projectOctoberPlan}. The older assembly provides a different
measured baseline:
\begin{align*}
\Delta P_{\mathrm{plan}} & = \SI{100}{\watt}-\SI{85}{\watt}
  = \SI{15}{\watt},\\
\Delta P_{\mathrm{record}} & = \SI{100}{\watt}-\SI{79.8928}{\watt}
  = \SI{20.1072}{\watt}\approx\SI{20.11}{\watt},\\
P_{\mathrm{in,logged}}-P_{\mathrm{out,calc}} & =
  \SI{104.418}{\watt}-\SI{79.8928}{\watt}
  = \SI{24.5252}{\watt}.
\end{align*}
The first two differences describe the earlier plan and V1 connector result; they do not identify loss mechanisms. Later V2 secondary-output tests delivered more than \SI{100}{\watt} with voltage droop and without thermal equilibrium. The third difference is logged input minus calculated output at the matched \SI{20}{\volt}, \SI{4}{\ampere} V1 point. Incomplete power-boundary and synchronization records prevent its allocation to component losses.

Secondary-winding heating and the V2 input hotspot suggest separate
magnetic and interconnect limitations. Winding current, resistance, and
temperature measurements are needed for the former; interlayer-path
inspection and voltage-drop measurements are needed for the suspected
single-via restriction. Neither has been isolated as the power limit.

The original \SI{85}{\watt} records should be recovered and repeated
with identified hardware and firmware. The \SI{79.8928}{\watt} result
belongs to the older assembly. Each hardware or MCU-port change requires
a comparison test with the board and settings recorded.

\begin{table}[htbp]
\centering
\caption[Measurements for diagnosing the power limit.]{Measurements needed before and after each change to diagnose the power limit.}
\label{tab:plan-gap-diagnostic-matrix}
\small
\begin{tabularx}{\textwidth}{@{}L{0.20\textwidth}L{0.26\textwidth}Y@{}}
\toprule
Question & Current result & Measurements needed\\
\midrule
Power readings and measurement points & Selected readings give the screening difference above; instruments disagree & Measure AC input, secondary-bus, connector, and auxiliary powers together, including cable drops and uncertainty. Calculate each stage loss from these readings.\\
Magnetic and interconnect loss & The installed magnetic path transfers power, but its individual losses are not measured & Identify the fitted assemblies and equivalent circuit; capture resonant current, winding/interconnect resistance, flux-related waveforms, and temperature. Use these to choose a magnetic, connection, or thermal change.\\
Semiconductor stress and commutation & Switching and ringing are visible; gate timing and calibrated commutation current are missing & Capture matched device voltage/current and gate timing, including burst starts and ends. Check peak/RMS stress and turn-on behavior before changing frequency, dead time, or drive settings.\\
Output stage and control limit & V1 VBUS remained regulated near its highest recorded load; later V2 secondary-output tests showed voltage droop; the active limiter is unidentified & Record secondary bus, VBUS, external load current, board-specific EN, commanded duty, and firmware state. Distinguish lost headroom, a duty limit, protection, and an output-stage limit.\\
Thermal and continuous operation & Labeled surface maxima rise at high load; dwell and hotspot identity are incomplete & Follow identified component temperatures, ambient, and airflow through the same defined-duration run. Check temperature margin at a steady load before increasing power or reducing volume.\\
Stability and protection & Load-removal response is retained without a complete timing or acceptance-band record & Repeat attachment, load increase/removal, PPS changes where supported, overload, and fault tests with event markers and declared bands/limits. Measure overshoot, recovery, shutdown, and restart.\\
\bottomrule
\end{tabularx}
\end{table}
\FloatBarrier

Controlled retests should vary one component or parameter at fixed input,
output setting, load, cooling, and initial temperature. Board and firmware,
changed values, synchronized power, waveforms, and temperatures should be
retained together. Lower temperature at one hotspot is useful only if regulation and all
other stress limits remain acceptable. Individual component losses have
not yet been separated.

Once the baseline meets its declared limits, the planned
\SI{90}{\watt}, \SI{95}{\watt}, and \SI{100}{\watt} steps can be
assessed with the same stress and thermal checks. A limiting result should
state the highest repeatable output, duration, and measured stop condition.
Sections~\ref{sec:revised-plan-criteria} and
\ref{sec:method-plan-validation} define the acceptance work.
